O objetivo é aplicar o Teorema $\Pi$ de Buckingham para confirmar o resultado do Problema 2.13, que afirma que a tensão $T$ numa corda de comprimento $L$, conectada a uma rocha de massa $m$ girando com velocidade $v$, satisfaz:
\[
\frac{TL}{mv^2} = \text{constante.}
\]

\textbf{Passo 1 — Identificação das variáveis e dimensões:}
\begin{table}[h]
\centering
\begin{tabular}{lll}
\toprule
\textbf{Grandeza} & \textbf{Símbolo} & \textbf{Dimensão} \\
\midrule
Tensão & $T$ & $M\,L\,T^{-2}$ \\
Comprimento da corda & $L$ & $L$ \\
Massa & $m$ & $M$ \\
Velocidade & $v$ & $L\,T^{-1}$ \\
\bottomrule
\end{tabular}
\end{table}

Temos $n = 4$ variáveis e $r = 3$ dimensões fundamentais $(M, L, T)$. Portanto, o número de grupos adimensionais é $k = n - r = 4 - 3 = 1$.

\textbf{Passo 2 — Escolha das variáveis base:}
Escolhemos $m$, $v$ e $L$ como variáveis base (contêm coletivamente $M$, $T$ e $L$, mas não formam grupo adimensional sozinhas).

\textbf{Passo 3 — Formação do grupo $\Pi_1$:}
\[
\Pi_1 = m^{a}\,v^{b}\,L^{c}\cdot T
\]
Substituindo as dimensões:
\[
M^0\,L^0\,T^0 = M^{a}\,(L\,T^{-1})^{b}\,L^{c}\cdot (M\,L\,T^{-2})
\]
\[
M^0\,L^0\,T^0 = M^{a+1}\;L^{b+c+1}\;T^{-b-2}
\]

\textbf{Passo 4 — Sistema de equações para os expoentes:}
\[
\begin{cases}
M:\; a + 1 = 0 &\Rightarrow a = -1 \\
T:\; -b - 2 = 0 &\Rightarrow b = -2 \\
L:\; b + c + 1 = 0 &\Rightarrow c = 1
\end{cases}
\]

\textbf{Passo 5 — Grupo adimensional resultante:}
\[
\Pi_1 = m^{-1}\,v^{-2}\,L^{1}\cdot T = \frac{TL}{mv^2}.
\]

\textbf{Conclusão:} O Teorema $\Pi$ de Buckingham confirma que $\Pi_1 = TL/(mv^2)$ é o único grupo adimensional independente, portanto $TL/(mv^2) = \text{constante}$, em acordo com o resultado do Problema 2.13.